\documentclass[11pt]{article}

\usepackage[margin=1in]{geometry}
\usepackage{amsmath}
\usepackage{amssymb}
\usepackage{xcolor}
\usepackage[most]{tcolorbox}
\usepackage{enumitem}
\usepackage{titlesec}
\usepackage{hyperref}

% Links map to the structural header blue.
% bookmarks=false avoids the spaced-jobname .out failure on Windows/tectonic.
\hypersetup{colorlinks=true, linkcolor=headerblue, urlcolor=headerblue, bookmarks=false}

\tcbuselibrary{skins}

% 1.1 Palette
\definecolor{headerblue}{HTML}{1C569E}
\definecolor{accentgreen}{HTML}{2E7D32}
\definecolor{accentorange}{HTML}{E27A1E}
\definecolor{workpurple}{HTML}{A020F0}
\definecolor{warnred}{HTML}{C0392B}

% 1.2 Subsection Typography
\titleformat{\subsection}
  {\normalfont\large\bfseries\color{headerblue}}{\thesubsection}{1em}{}

% 2. Table of Contents Depth
\setcounter{secnumdepth}{2}
\setcounter{tocdepth}{2}

% 3. Box System (Unbreakable)
\newtcolorbox{topicsbox}{
    enhanced,
    colback=headerblue!6!white, colframe=headerblue,
    coltitle=white, fonttitle=\bfseries,
    title={Evaluating Integrals \& Integral Equations at a Glance: Core Sub-Topics},
    boxrule=0.6pt, arc=2mm,
    left=3mm, right=3mm, top=2mm, bottom=2mm
}

\newtcolorbox{formulabox}[1][Master Formula]{
    enhanced,
    colback=accentorange!8!white, colframe=accentorange,
    coltitle=white, fonttitle=\bfseries,
    title={#1},
    boxrule=0.6pt, arc=1mm,
    left=3mm, right=3mm, top=2mm, bottom=2mm
}

\newtcolorbox{methodbox}[1][Method]{
    enhanced,
    colback=accentgreen!6!white, colframe=accentgreen,
    coltitle=white, fonttitle=\bfseries,
    title={#1},
    boxrule=0.6pt, arc=1mm,
    left=3mm, right=3mm, top=2mm, bottom=2mm
}

\newtcolorbox{examplebox}[1][Worked Example]{
    enhanced,
    colback=workpurple!4!white, colframe=workpurple,
    coltitle=white, fonttitle=\bfseries,
    title={#1},
    boxrule=0.4pt, sharp corners,
    borderline west={3pt}{0pt}{workpurple},
    left=5mm, right=3mm, top=2mm, bottom=2mm
}

\newtcolorbox{conceptbox}[1][Key Takeaway]{
    enhanced,
    colback=accentorange!5!white, colframe=accentorange,
    coltitle=white, fonttitle=\bfseries,
    title={#1},
    boxrule=0.3pt, arc=1mm,
    left=3mm, right=3mm, top=2mm, bottom=2mm
}

\newtcolorbox{warningbox}[1][Common Pitfall]{
    enhanced,
    colback=red!3!white, colframe=warnred,
    coltitle=white, fonttitle=\bfseries,
    title={#1},
    boxrule=0.4pt, sharp corners,
    borderline west={3pt}{0pt}{warnred},
    left=5mm, right=3mm, top=2mm, bottom=2mm
}

\begin{document}

\begin{center}
    {\Large\bfseries\color{headerblue} Evaluating Integrals\\[2pt] and Solving Integral Equations with Laplace Transforms}\\[6pt]
    \large A Self-Study Reference \& Master Guide --- Sections 2.7--2.8
\end{center}
\vspace{0.4cm}

\tableofcontents
\vspace{0.8cm}

%======================================================================
\phantomsection
\addcontentsline{toc}{section}{Part I --- Theory \& Worked Examples}
\section*{Part I --- Theory \& Worked Examples}
\setcounter{section}{2}
\setcounter{subsection}{6}

\begin{topicsbox}
\begin{itemize}[leftmargin=*, itemsep=2pt]
  \item \textbf{The Laplace Method for Integrals (2.7):} a definite improper integral $\int_0^\infty f(t)\,dt$ is nothing more than the Laplace transform $F(s)$ evaluated at $s=0$, because setting $s=0$ erases the kernel $e^{-st}$ and leaves the bare integral behind. Paired with the \emph{division-by-$t$} property, this turns hard Calculus~II integrals --- the kind that have no elementary antiderivative --- into one-line transform lookups.
  \item \textbf{The Laplace Method for Integral Equations (2.8):} an equation that hides its unknown $x(t)$ \emph{inside} an integral of the form $\int_0^t K(t-\tau)\,x(\tau)\,d\tau$ is a \emph{convolution} in disguise. The Convolution Theorem converts that integral into the ordinary product $K(s)\,X(s)$, collapsing the whole equation into linear algebra for $X(s)$.
  \item \textbf{One toolkit, two directions:} both methods rest on the same machinery --- transform, exploit a structural property ($F(0)$, integration-in-$s$, or the convolution product), then either read off a number or invert back to $x(t)$. The Laplace transform is a bridge that carries calculus problems into algebra and walks the answers back.
\end{itemize}
\end{topicsbox}

\begin{conceptbox}[Why a Transform Can Evaluate an Integral You Cannot Antidifferentiate]
The deepest idea in this section is almost a sleight of hand. The definition of the transform,
\[ F(s) = \int_0^\infty e^{-st} f(t)\,dt, \]
already \emph{contains} an integral from $0$ to $\infty$. The factor $e^{-st}$ is the only thing standing between $F(s)$ and the raw integral $\int_0^\infty f(t)\,dt$. \textcolor{accentorange}{Set $s=0$ and that factor becomes $e^{0}=1$}, so the kernel vanishes and $F(0)$ \emph{is} the integral we wanted. This is why a problem that looks like Calculus~II --- an improper integral with no elementary antiderivative, such as $\int_0^\infty \frac{\sin t}{t}\,dt$ --- can be cracked with a transform table instead of an antiderivative. We never antidifferentiate $\frac{\sin t}{t}$ (it has no elementary antiderivative); we transform it and evaluate at a point.
\end{conceptbox}

\subsection{Evaluating Improper Integrals (Calculus II Style)}

This first method answers a question that haunts every Calculus~II student: \emph{what do you do with an improper integral that has no elementary antiderivative?} The Laplace transform offers an escape hatch. Two properties do all the work --- the evaluation-at-$s=0$ identity, and the division-by-$t$ (integration-in-$s$) property --- and together they evaluate integrals that ordinary calculus cannot touch.

\begin{formulabox}[Property A --- Evaluation at $s=0$]
If $\mathcal{L}\{f(t)\} = F(s)$ and the improper integral converges, then
\[ \int_0^\infty f(t)\,dt = F(0). \]
The proof is immediate from the definition: $F(s) = \int_0^\infty e^{-st} f(t)\,dt$, so substituting $s=0$ replaces $e^{-st}$ by $e^{0}=1$, leaving $F(0) = \int_0^\infty f(t)\,dt$ exactly. The transform variable $s$ is a ``dial,'' and turning it to zero recovers the undamped integral.
\end{formulabox}

\begin{formulabox}[Property B --- Division by $t$ (Integration in the $s$-Domain)]
If $\mathcal{L}\{f(t)\} = F(s)$, then dividing the time function by $t$ corresponds to \emph{integrating} the transform from $s$ out to infinity:
\[ \mathcal{L}\!\left\{\frac{f(t)}{t}\right\} = \int_s^\infty F(\sigma)\,d\sigma. \]
This is the exact mirror of the familiar rule $\mathcal{L}\{t\,f(t)\} = -F'(s)$: multiplying by $t$ \emph{differentiates} the transform (with a sign), so dividing by $t$ \emph{integrates} it. The dummy variable $\sigma$ sweeps from the current $s$ out to $\infty$; the surviving variable is $s$.
\end{formulabox}

\noindent Why do we need Property~B at all? Because the integrals worth evaluating this way --- like $\int_0^\infty \frac{\sin t}{t}\,dt$ --- carry a stubborn $\frac{1}{t}$ that no table entry handles directly. Property~B is the tool that absorbs that $\frac{1}{t}$: it tells us the transform of $\frac{f(t)}{t}$ without ever performing the division in the time domain. Once we have $\mathcal{L}\{f(t)/t\}$ as a clean function of $s$, Property~A finishes the job by evaluating it at the right point.

\begin{methodbox}[The Recipe --- Evaluating $\int_0^\infty \tfrac{f(t)}{t}\,e^{-at}\,dt$]
\textbf{Step 1 --- Strip the integral down to a transform.} Recognize that $\int_0^\infty e^{-at}\,\dfrac{f(t)}{t}\,dt$ is simply $\mathcal{L}\!\left\{\dfrac{f(t)}{t}\right\}$ evaluated at $s=a$. (When there is no exponential, $a=0$ and you evaluate the limit as $s\to0^+$.)

\smallskip
\textbf{Step 2 --- Find the base transform.} Look up $F(s) = \mathcal{L}\{f(t)\}$ from the table.

\smallskip
\textbf{Step 3 --- Apply division-by-$t$.} Compute $\displaystyle G(s) = \mathcal{L}\!\left\{\frac{f(t)}{t}\right\} = \int_s^\infty F(\sigma)\,d\sigma$. This integral in $\sigma$ is the heart of the method.

\smallskip
\textbf{Step 4 --- Evaluate at the right point.} Plug in $s=a$ to read off the numerical value: the answer is $G(a)$. If $a=0$, take the limit $\displaystyle\lim_{s\to0^+}G(s)$.
\end{methodbox}

\begin{examplebox}[Worked Example 1: $\displaystyle \int_0^\infty \frac{e^{-t}\sin t}{t}\,dt$]
Evaluate the improper integral $\displaystyle \int_0^\infty \frac{e^{-t}\sin t}{t}\,dt$, which has no elementary antiderivative.

\smallskip
\textbf{Step 1 --- Read the integral as a transform evaluated at a point.} The factor $e^{-t}$ is precisely the Laplace kernel $e^{-st}$ with $s=1$. So if we let $f(t) = \sin t$, the integral is the transform of $\dfrac{\sin t}{t}$ evaluated at $s=1$:
\[ \color{workpurple} \int_0^\infty \frac{e^{-t}\sin t}{t}\,dt = \int_0^\infty e^{-st}\,\frac{\sin t}{t}\,dt\,\Big|_{s=1} = \mathcal{L}\!\left\{\frac{\sin t}{t}\right\}\Big|_{s=1}. \]

\textbf{Step 2 --- Base transform.} From the table, $\mathcal{L}\{\sin t\} = \dfrac{1}{s^2+1}$, so $F(\sigma) = \dfrac{1}{\sigma^2+1}$.

\smallskip
\textbf{Step 3 --- Apply division-by-$t$.} By Property~B,
\[ \color{workpurple} G(s) = \mathcal{L}\!\left\{\frac{\sin t}{t}\right\} = \int_s^\infty \frac{1}{\sigma^2+1}\,d\sigma. \]
The antiderivative of $\dfrac{1}{\sigma^2+1}$ is $\arctan\sigma$, so we evaluate the improper integral as a limit:
\[ \color{workpurple} G(s) = \Big[\arctan\sigma\Big]_{\sigma=s}^{\sigma\to\infty} = \lim_{\sigma\to\infty}\arctan\sigma - \arctan s = \frac{\pi}{2} - \arctan s. \]
The combination $\frac{\pi}{2} - \arctan s$ is exactly $\operatorname{arccot} s = \arctan\!\frac{1}{s}$, so equivalently $G(s) = \arctan\!\dfrac{1}{s}$.

\smallskip
\textbf{Step 4 --- Evaluate at $s=1$.} Set $s=1$ in $G(s)$:
\[ \color{workpurple} \int_0^\infty \frac{e^{-t}\sin t}{t}\,dt = G(1) = \frac{\pi}{2} - \arctan 1 = \frac{\pi}{2} - \frac{\pi}{4} = \frac{\pi}{4}. \]
The whole integral collapses to $\boxed{\dfrac{\pi}{4}}$ --- found without ever antidifferentiating $\frac{\sin t}{t}$.
\end{examplebox}

\begin{conceptbox}[The Limit as $s\to0$: A Free Bonus]
Notice what $G(s) = \dfrac{\pi}{2} - \arctan s$ gives if we instead let $s \to 0^+$ --- that is, if we drop the damping factor $e^{-t}$ entirely and ask for the \emph{undamped} integral:
\[ \int_0^\infty \frac{\sin t}{t}\,dt = \lim_{s\to0^+}\left(\frac{\pi}{2} - \arctan s\right) = \frac{\pi}{2} - 0 = \frac{\pi}{2}. \]
This is the celebrated \textcolor{accentorange}{Dirichlet integral}. The same $G(s)$ that gave us $\pi/4$ at $s=1$ delivers the famous $\pi/2$ in the limit $s\to0$ --- a single transform answers a whole family of integrals, one for each value of the dial $s$. You will reprove this directly in Problem~1.
\end{conceptbox}

\subsection{Solving Volterra Integral Equations}

We now turn from \emph{evaluating} integrals to \emph{solving for an unknown trapped inside one}. A \textbf{Volterra integral equation of convolution type} has the form
\[ x(t) = f(t) + \int_0^t K(t-\tau)\,x(\tau)\,d\tau, \]
where $f(t)$ (the forcing) and $K(t)$ (the \emph{kernel}) are known, and $x(t)$ is the unknown we must find. The decisive observation is that the integral is a \textcolor{accentorange}{convolution}: $\int_0^t K(t-\tau)\,x(\tau)\,d\tau = (K * x)(t)$. That single recognition unlocks everything.

\begin{formulabox}[The Convolution Theorem --- The Master Key]
The transform of a convolution is the \emph{product} of the transforms:
\[ \mathcal{L}\{(K * x)(t)\} = \mathcal{L}\!\left\{\int_0^t K(t-\tau)\,x(\tau)\,d\tau\right\} = K(s)\,X(s). \]
An operation that is genuinely hard in the time domain (the convolution integral) becomes ordinary multiplication in the $s$-domain. This is why integral equations succumb to Laplace methods: the integral that defines the equation simply turns into $K(s)\,X(s)$.
\end{formulabox}

\begin{methodbox}[The Recipe --- Solving the Convolution Equation]
\textbf{Step 1 --- Transform every term.} Apply $\mathcal{L}$ to $x(t) = f(t) + (K*x)(t)$. By the Convolution Theorem the integral becomes $K(s)X(s)$:
\[ X(s) = F(s) + K(s)\,X(s). \]

\smallskip
\textbf{Step 2 --- Gather the $X(s)$ terms.} Move $K(s)X(s)$ to the left:
\[ X(s) - K(s)\,X(s) = F(s). \]

\smallskip
\textbf{Step 3 --- Factor out $X(s)$.}
\[ X(s)\,\big(1 - K(s)\big) = F(s). \]

\smallskip
\textbf{Step 4 --- Isolate $X(s)$} by dividing by the bracket:
\[ X(s) = \frac{F(s)}{1 - K(s)}. \]

\smallskip
\textbf{Step 5 --- Decompose and invert} via partial fractions to recover $x(t) = \mathcal{L}^{-1}\{X(s)\}$.
\end{methodbox}

\begin{examplebox}[Worked Example 2: $\displaystyle x(t) = t^2 + \int_0^t \sin(t-\tau)\,x(\tau)\,d\tau$]
Solve the Volterra equation $\displaystyle x(t) = t^{2} + \int_0^t \sin(t-\tau)\,x(\tau)\,d\tau$.

\smallskip
\textbf{Step 1 --- Identify the pieces.} The forcing is $f(t) = t^2$ and the kernel is $K(t) = \sin t$. From the table,
\[ \color{workpurple} F(s) = \mathcal{L}\{t^2\} = \frac{2!}{s^3} = \frac{2}{s^3}, \qquad K(s) = \mathcal{L}\{\sin t\} = \frac{1}{s^2+1}. \]

\textbf{Step 2 --- Transform the equation.} The integral is the convolution $(\sin * x)(t)$, which transforms to $K(s)X(s)$:
\[ \color{workpurple} X(s) = \frac{2}{s^3} + \frac{1}{s^2+1}\,X(s). \]

\textbf{Step 3 --- Gather $X(s)$ on the left and factor.}
\[ \color{workpurple} X(s) - \frac{1}{s^2+1}\,X(s) = \frac{2}{s^3} \quad\Longrightarrow\quad X(s)\left(1 - \frac{1}{s^2+1}\right) = \frac{2}{s^3}. \]

\textbf{Step 4 --- Combine the bracket over a common denominator.} Write $1 = \dfrac{s^2+1}{s^2+1}$ so the two fractions share the denominator $s^2+1$, then subtract:
\[ \color{workpurple} 1 - \frac{1}{s^2+1} = \frac{s^2+1}{s^2+1} - \frac{1}{s^2+1} = \frac{(s^2+1) - 1}{s^2+1} = \frac{s^2}{s^2+1}. \]
So the equation becomes
\[ \color{workpurple} X(s)\left(\frac{s^2}{s^2+1}\right) = \frac{2}{s^3}. \]

\textbf{Step 5 --- Multiply by the reciprocal to isolate $X(s)$.} The reciprocal of $\dfrac{s^2}{s^2+1}$ is $\dfrac{s^2+1}{s^2}$:
\[ \color{workpurple} X(s) = \frac{2}{s^3}\cdot\frac{s^2+1}{s^2} = \frac{2(s^2+1)}{s^{5}}. \]
The $s$-domain algebra is finished. The decomposition and inversion follow in the box below.
\end{examplebox}

\begin{examplebox}[Worked Example 2 (continued): Decomposition and Inversion]
We carry $X(s) = \dfrac{2(s^2+1)}{s^{5}}$ forward.

\smallskip
\textbf{Step 6 --- Set up the partial fraction decomposition.} The denominator is the single repeated factor $s^{5}$, so a complete decomposition assigns one term to each power from $s^{1}$ up to $s^{5}$:
\[ \color{workpurple} \frac{2(s^2+1)}{s^{5}} = \frac{A}{s} + \frac{B}{s^2} + \frac{C}{s^3} + \frac{D}{s^4} + \frac{E}{s^5}. \]
Multiply both sides by $s^{5}$ to clear every denominator:
\[ \color{workpurple} 2(s^2+1) = A\,s^4 + B\,s^3 + C\,s^2 + D\,s + E. \]

\textbf{Step 7 --- Match coefficients power by power.} Expand the left side as $2s^2 + 2$, i.e.\ $0\cdot s^4 + 0\cdot s^3 + 2\cdot s^2 + 0\cdot s + 2$, and equate each coefficient:
\[ \color{workpurple} \underbrace{A = 0}_{s^4}, \quad \underbrace{B = 0}_{s^3}, \quad \underbrace{C = 2}_{s^2}, \quad \underbrace{D = 0}_{s^1}, \quad \underbrace{E = 2}_{s^0}. \]
The formal decomposition therefore collapses to just two surviving terms:
\[ \color{workpurple} X(s) = \frac{2}{s^3} + \frac{2}{s^5}. \]
(This matches what the eye sees if it splits the numerator directly: $\dfrac{2(s^2+1)}{s^5} = \dfrac{2s^2}{s^5} + \dfrac{2}{s^5} = \dfrac{2}{s^3} + \dfrac{2}{s^5}$. The full setup confirms there are no hidden $1/s$, $1/s^2$, or $1/s^4$ pieces.)

\smallskip
\textbf{Step 8 --- Invert term by term.} Use $\mathcal{L}^{-1}\!\left\{\dfrac{1}{s^{n+1}}\right\} = \dfrac{t^{n}}{n!}$. For the first term, $\dfrac{2}{s^3} = 2\cdot\dfrac{1}{s^3} \to 2\cdot\dfrac{t^2}{2!} = 2\cdot\dfrac{t^2}{2} = t^2$. For the second, $\dfrac{2}{s^5} = 2\cdot\dfrac{1}{s^5} \to 2\cdot\dfrac{t^4}{4!} = 2\cdot\dfrac{t^4}{24} = \dfrac{t^4}{12}$:
\[ \color{workpurple} x(t) = t^2 + \frac{t^4}{12}. \]

\textbf{Step 9 --- Verify.} The forcing $t^2$ reappears as the leading term, so the convolution must supply exactly the $\dfrac{t^4}{12}$ correction. Indeed $\mathcal{L}\!\left\{\dfrac{t^4}{12}\right\} = \dfrac{1}{12}\cdot\dfrac{4!}{s^5} = \dfrac{2}{s^5}$, while $K(s)X(s) = \dfrac{1}{s^2+1}\cdot\dfrac{2(s^2+1)}{s^5} = \dfrac{2}{s^5}$ --- the two agree, so $x(t) = t^2 + \frac{t^4}{12}$ solves the equation. \qquad$\checkmark$
\end{examplebox}

\begin{warningbox}[Mind the Argument: $K(t-\tau)$, Not $K(\tau)$]
The Convolution Theorem only applies when the kernel's argument is the \emph{difference} $t-\tau$. An equation with $\int_0^t K(t-\tau)\,x(\tau)\,d\tau$ is a true convolution and transforms to $K(s)X(s)$. But $\int_0^t K(t)\,x(\tau)\,d\tau$ (kernel depending on $t$ alone, pulled outside) or a kernel written as $K(t+\tau)$ is \emph{not} a convolution --- the product rule $K(s)X(s)$ does \textbf{not} apply. Always confirm the integrand has the clean $t-\tau$ structure before invoking the theorem.
\end{warningbox}

\begin{conceptbox}[Key Takeaway]
Both halves of this section run on the transform's structure. To \textcolor{accentorange}{evaluate} an improper integral, recognize it as $F(s)$ at a point --- $s=0$ for the raw integral, or $s=a$ when an $e^{-at}$ rides along --- and lean on \textcolor{accentorange}{division-by-$t$} to absorb a $\frac1t$. To \textcolor{accentorange}{solve} an integral equation, recognize the integral as a \textcolor{accentorange}{convolution}, replace it with the product $K(s)X(s)$, and solve the one-line algebra $X(s) = \dfrac{F(s)}{1-K(s)}$.
\end{conceptbox}

%======================================================================
\clearpage
\phantomsection
\addcontentsline{toc}{section}{Part II --- Practice Set}
\section*{Part II --- Practice Set}

Work each problem before turning to the complete solutions in Part~III. Problems 1--4 evaluate improper integrals, Problems 5--8 solve Volterra integral equations (Problem~8 is integro-differential), and Problems 9--10 ask you to prove the two properties that power this section.

\bigskip
\noindent\textbf{Problem 1: The Dirichlet Integral.}\\[2pt]
Evaluate $\displaystyle \int_0^\infty \frac{\sin t}{t}\,dt$ by finding $\mathcal{L}\{\sin t/t\}$ and taking the limit as $s\to0^+$.

\bigskip
\noindent\textbf{Problem 2: A Frullani-Type Integral.}\\[2pt]
Evaluate $\displaystyle \int_0^\infty \frac{e^{-t} - e^{-3t}}{t}\,dt$.

\bigskip
\noindent\textbf{Problem 3: Evaluation at a Point.}\\[2pt]
Evaluate $\displaystyle \int_0^\infty t\,e^{-2t}\cos t\,dt$ by recognizing it as a transform evaluated at $s=2$.

\bigskip
\noindent\textbf{Problem 4: Division-by-$t$ with Damping.}\\[2pt]
Evaluate $\displaystyle \int_0^\infty \frac{e^{-2t}\sin(3t)}{t}\,dt$.

\bigskip
\noindent\textbf{Problem 5: A Volterra Equation.}\\[2pt]
Solve for $x(t)$:
\[ x(t) = 1 + \int_0^t (t-\tau)\,x(\tau)\,d\tau. \]

\bigskip
\noindent\textbf{Problem 6: A Volterra Equation.}\\[2pt]
Solve for $x(t)$:
\[ x(t) = e^{t} + \int_0^t e^{\,t-\tau}\,x(\tau)\,d\tau. \]

\bigskip
\noindent\textbf{Problem 7: A Volterra Equation.}\\[2pt]
Solve for $x(t)$:
\[ x(t) = t + \int_0^t \sin(t-\tau)\,x(\tau)\,d\tau. \]

\bigskip
\noindent\textbf{Problem 8: An Integro-Differential Equation.}\\[2pt]
Solve for $x(t)$, given $x(0) = 0$:
\[ x'(t) = 1 - \int_0^t x(\tau)\,d\tau. \]

\bigskip
\noindent\textbf{Problem 9 (Proof): The Integral Property of Laplace.}\\[2pt]
Prove the division-by-$t$ property
\[ \mathcal{L}\!\left\{\frac{f(t)}{t}\right\} = \int_s^\infty F(\sigma)\,d\sigma, \]
and deduce the evaluation corollary $\displaystyle \int_0^\infty \frac{f(t)}{t}\,dt = \int_0^\infty F(\sigma)\,d\sigma$.

\bigskip
\noindent\textbf{Problem 10 (Proof): Inversion of the Convolution Product.}\\[2pt]
Prove the Convolution Theorem --- that the product $F(s)\,G(s)$ inverts to the convolution:
\[ \mathcal{L}\{(f * g)(t)\} = F(s)\,G(s), \qquad (f*g)(t) = \int_0^t f(\tau)\,g(t-\tau)\,d\tau. \]

%======================================================================
\clearpage
\phantomsection
\addcontentsline{toc}{section}{Part III --- Complete Solutions}
\section*{Part III --- Complete Solutions}

\noindent\textbf{Solution 1: The Dirichlet Integral $\displaystyle \int_0^\infty \frac{\sin t}{t}\,dt$.}\\[2pt]
\textbf{Step 1 --- Transform $\sin t / t$ via division-by-$t$.} With $f(t) = \sin t$ and $F(\sigma) = \dfrac{1}{\sigma^2+1}$, Property~B gives
\[ \color{workpurple} \mathcal{L}\!\left\{\frac{\sin t}{t}\right\} = \int_s^\infty \frac{1}{\sigma^2+1}\,d\sigma = \Big[\arctan\sigma\Big]_s^\infty = \frac{\pi}{2} - \arctan s. \]
\textbf{Step 2 --- Recognize the raw integral as the $s\to0^+$ limit.} The undamped integral is $\mathcal{L}\{\sin t/t\}$ at $s=0$:
\[ \color{workpurple} \int_0^\infty \frac{\sin t}{t}\,dt = \lim_{s\to0^+}\left(\frac{\pi}{2} - \arctan s\right) = \frac{\pi}{2} - 0 = \frac{\pi}{2}. \]
The Dirichlet integral equals $\dfrac{\pi}{2}$, obtained with no antiderivative of $\sin t / t$. \checkmark

\bigskip
\noindent\textbf{Solution 2: $\displaystyle \int_0^\infty \frac{e^{-t} - e^{-3t}}{t}\,dt$.}\\[2pt]
\textbf{Step 1 --- Base transform.} With $f(t) = e^{-t} - e^{-3t}$,
\[ \color{workpurple} F(\sigma) = \frac{1}{\sigma+1} - \frac{1}{\sigma+3}. \]
\textbf{Step 2 --- Apply division-by-$t$.} Integrate $F$ from $s$ to $\infty$. The antiderivative is $\ln(\sigma+1) - \ln(\sigma+3) = \ln\dfrac{\sigma+1}{\sigma+3}$:
\[ \color{workpurple} \mathcal{L}\!\left\{\frac{f(t)}{t}\right\} = \int_s^\infty\!\left(\frac{1}{\sigma+1} - \frac{1}{\sigma+3}\right)d\sigma = \left[\ln\frac{\sigma+1}{\sigma+3}\right]_s^\infty. \]
As $\sigma\to\infty$, $\dfrac{\sigma+1}{\sigma+3}\to1$ and $\ln1 = 0$, so
\[ \color{workpurple} \mathcal{L}\!\left\{\frac{f(t)}{t}\right\} = 0 - \ln\frac{s+1}{s+3} = \ln\frac{s+3}{s+1}. \]
\textbf{Step 3 --- Evaluate at $s=0$.} The integral has no exponential damping ($a=0$), so set $s=0$:
\[ \color{workpurple} \int_0^\infty \frac{e^{-t} - e^{-3t}}{t}\,dt = \ln\frac{0+3}{0+1} = \ln 3. \]
(This is the Frullani value $\ln(b/a)$ with $a=1$, $b=3$.) \checkmark

\bigskip
\noindent\textbf{Solution 3: $\displaystyle \int_0^\infty t\,e^{-2t}\cos t\,dt$.}\\[2pt]
\textbf{Step 1 --- Recognize the transform.} The integral is $\mathcal{L}\{t\cos t\}$ evaluated at $s=2$, since the $e^{-2t}$ is the kernel with $s=2$:
\[ \color{workpurple} \int_0^\infty t\,e^{-2t}\cos t\,dt = \mathcal{L}\{t\cos t\}\big|_{s=2}. \]
\textbf{Step 2 --- Find $\mathcal{L}\{t\cos t\}$ by the multiply-by-$t$ rule.} Using $\mathcal{L}\{t\,f(t)\} = -F'(s)$ with $\mathcal{L}\{\cos t\} = \dfrac{s}{s^2+1}$:
\[ \color{workpurple} \mathcal{L}\{t\cos t\} = -\frac{d}{ds}\!\left(\frac{s}{s^2+1}\right) = -\frac{(1)(s^2+1) - s(2s)}{(s^2+1)^2} = -\frac{1 - s^2}{(s^2+1)^2} = \frac{s^2-1}{(s^2+1)^2}. \]
\textbf{Step 3 --- Evaluate at $s=2$.}
\[ \color{workpurple} \int_0^\infty t\,e^{-2t}\cos t\,dt = \frac{2^2 - 1}{(2^2+1)^2} = \frac{3}{25}. \]
\checkmark

\bigskip
\noindent\textbf{Solution 4: $\displaystyle \int_0^\infty \frac{e^{-2t}\sin(3t)}{t}\,dt$.}\\[2pt]
\textbf{Step 1 --- Read as a transform at $s=2$.} With $f(t) = \sin 3t$, the integral is $\mathcal{L}\{\sin 3t / t\}$ at $s=2$.

\smallskip
\textbf{Step 2 --- Division-by-$t$.} Here $F(\sigma) = \dfrac{3}{\sigma^2+9}$, whose antiderivative is $\arctan\dfrac{\sigma}{3}$:
\[ \color{workpurple} \mathcal{L}\!\left\{\frac{\sin 3t}{t}\right\} = \int_s^\infty \frac{3}{\sigma^2+9}\,d\sigma = \left[\arctan\frac{\sigma}{3}\right]_s^\infty = \frac{\pi}{2} - \arctan\frac{s}{3}. \]
\textbf{Step 3 --- Evaluate at $s=2$.} Using $\dfrac{\pi}{2} - \arctan x = \arctan\dfrac{1}{x}$ with $x = \dfrac{2}{3}$:
\[ \color{workpurple} \int_0^\infty \frac{e^{-2t}\sin 3t}{t}\,dt = \frac{\pi}{2} - \arctan\frac{2}{3} = \arctan\frac{3}{2} \approx 0.9828. \]
\checkmark

\bigskip
\noindent\textbf{Solution 5: $\displaystyle x(t) = 1 + \int_0^t (t-\tau)\,x(\tau)\,d\tau$.}\\[2pt]
\textbf{Step 1 --- Identify and transform.} The kernel is $K(t) = t$ with $K(s) = \dfrac{1}{s^2}$, and $f(t) = 1$ with $F(s) = \dfrac{1}{s}$. The integral is the convolution $(t * x)(t) \to K(s)X(s)$:
\[ \color{workpurple} X(s) = \frac{1}{s} + \frac{1}{s^2}\,X(s). \]
\textbf{Step 2 --- Gather and factor.}
\[ \color{workpurple} X(s) - \frac{1}{s^2}\,X(s) = \frac{1}{s} \quad\Longrightarrow\quad X(s)\left(1 - \frac{1}{s^2}\right) = \frac{1}{s}. \]
\textbf{Step 3 --- Common denominator in the bracket.}
\[ \color{workpurple} 1 - \frac{1}{s^2} = \frac{s^2 - 1}{s^2}, \qquad\text{so}\qquad X(s)\left(\frac{s^2-1}{s^2}\right) = \frac{1}{s}. \]
\textbf{Step 4 --- Multiply by the reciprocal and cancel.}
\[ \color{workpurple} X(s) = \frac{1}{s}\cdot\frac{s^2}{s^2-1} = \frac{s}{s^2-1}. \]
\textbf{Step 5 --- Invert.} Recognize $\dfrac{s}{s^2-1}$ as the transform of $\cosh t$:
\[ \color{workpurple} x(t) = \cosh t. \]
\noindent\textit{Check.} $1 + \int_0^t (t-\tau)\cosh\tau\,d\tau = 1 + (\cosh t - 1) = \cosh t = x(t)$. \checkmark

\bigskip
\noindent\textbf{Solution 6: $\displaystyle x(t) = e^{t} + \int_0^t e^{\,t-\tau}\,x(\tau)\,d\tau$.}\\[2pt]
\textbf{Step 1 --- Identify and transform.} Here $K(t) = e^{t}$ with $K(s) = \dfrac{1}{s-1}$, and $f(t) = e^t$ with $F(s) = \dfrac{1}{s-1}$. The integral is $(e^t * x)(t) \to K(s)X(s)$:
\[ \color{workpurple} X(s) = \frac{1}{s-1} + \frac{1}{s-1}\,X(s). \]
\textbf{Step 2 --- Gather and factor.}
\[ \color{workpurple} X(s)\left(1 - \frac{1}{s-1}\right) = \frac{1}{s-1}. \]
\textbf{Step 3 --- Common denominator in the bracket.} Write $1 = \dfrac{s-1}{s-1}$:
\[ \color{workpurple} 1 - \frac{1}{s-1} = \frac{(s-1) - 1}{s-1} = \frac{s-2}{s-1}, \qquad\text{so}\qquad X(s)\left(\frac{s-2}{s-1}\right) = \frac{1}{s-1}. \]
\textbf{Step 4 --- Multiply by the reciprocal and cancel.} The factors of $(s-1)$ cancel:
\[ \color{workpurple} X(s) = \frac{1}{s-1}\cdot\frac{s-1}{s-2} = \frac{1}{s-2}. \]
\textbf{Step 5 --- Invert.}
\[ \color{workpurple} x(t) = e^{2t}. \]
\noindent\textit{Check.} $e^t + \int_0^t e^{t-\tau}e^{2\tau}\,d\tau = e^t + e^t\!\int_0^t e^{\tau}\,d\tau = e^t + e^t(e^t - 1) = e^{2t} = x(t)$. \checkmark

\bigskip
\noindent\textbf{Solution 7: $\displaystyle x(t) = t + \int_0^t \sin(t-\tau)\,x(\tau)\,d\tau$.}\\[2pt]
\textbf{Step 1 --- Identify and transform.} Kernel $K(t) = \sin t$ with $K(s) = \dfrac{1}{s^2+1}$; forcing $f(t) = t$ with $F(s) = \dfrac{1}{s^2}$:
\[ \color{workpurple} X(s) = \frac{1}{s^2} + \frac{1}{s^2+1}\,X(s). \]
\textbf{Step 2 --- Gather, factor, and combine the bracket.}
\[ \color{workpurple} X(s)\left(1 - \frac{1}{s^2+1}\right) = \frac{1}{s^2}, \qquad 1 - \frac{1}{s^2+1} = \frac{s^2}{s^2+1}, \qquad\text{so}\qquad X(s)\,\frac{s^2}{s^2+1} = \frac{1}{s^2}. \]
\textbf{Step 3 --- Multiply by the reciprocal.}
\[ \color{workpurple} X(s) = \frac{1}{s^2}\cdot\frac{s^2+1}{s^2} = \frac{s^2+1}{s^4} = \frac{1}{s^2} + \frac{1}{s^4}. \]
\textbf{Step 4 --- Invert term by term.} Using $\mathcal{L}^{-1}\{1/s^2\} = t$ and $\mathcal{L}^{-1}\{1/s^4\} = \dfrac{t^3}{3!} = \dfrac{t^3}{6}$:
\[ \color{workpurple} x(t) = t + \frac{t^3}{6}. \]
\noindent\textit{Check.} The forcing $t$ reappears, and the convolution must supply $\dfrac{t^3}{6}$; indeed $K(s)X(s) = \dfrac{1}{s^2+1}\cdot\dfrac{s^2+1}{s^4} = \dfrac{1}{s^4} \to \dfrac{t^3}{6}$. \checkmark

\bigskip
\noindent\textbf{Solution 8: $\displaystyle x'(t) = 1 - \int_0^t x(\tau)\,d\tau, \quad x(0) = 0$.}\\[2pt]
\textbf{Step 1 --- Transform each term.} The derivative rule gives $\mathcal{L}\{x'\} = sX(s) - x(0) = sX(s)$ since $x(0)=0$. The accumulation integral $\int_0^t x(\tau)\,d\tau$ is the convolution of $x$ with the constant $1$, so it transforms to $\dfrac{X(s)}{s}$; and $\mathcal{L}\{1\} = \dfrac{1}{s}$:
\[ \color{workpurple} sX(s) = \frac{1}{s} - \frac{X(s)}{s}. \]
\textbf{Step 2 --- Gather the $X(s)$ terms.} Add $\dfrac{X(s)}{s}$ to both sides:
\[ \color{workpurple} sX(s) + \frac{X(s)}{s} = \frac{1}{s}. \]
\textbf{Step 3 --- Factor out $X(s)$.}
\[ \color{workpurple} X(s)\left(s + \frac{1}{s}\right) = \frac{1}{s}. \]
\textbf{Step 4 --- Common denominator in the bracket.}
\[ \color{workpurple} s + \frac{1}{s} = \frac{s^2 + 1}{s}, \qquad\text{so}\qquad X(s)\left(\frac{s^2+1}{s}\right) = \frac{1}{s}. \]
\textbf{Step 5 --- Multiply by the reciprocal and cancel.} The factors of $s$ cancel:
\[ \color{workpurple} X(s) = \frac{1}{s}\cdot\frac{s}{s^2+1} = \frac{1}{s^2+1}. \]
\textbf{Step 6 --- Invert.}
\[ \color{workpurple} x(t) = \sin t. \]
\noindent\textit{Check.} $x(0) = \sin 0 = 0$ \checkmark; and $x'(t) = \cos t$, while $1 - \int_0^t \sin\tau\,d\tau = 1 - (1 - \cos t) = \cos t$. \checkmark

\bigskip
\noindent\textbf{Solution 9 (Proof): The Integral Property of Laplace.}\\[2pt]
\textbf{Goal.} Show $\displaystyle \mathcal{L}\!\left\{\frac{f(t)}{t}\right\} = \int_s^\infty F(\sigma)\,d\sigma$.

\smallskip
\textbf{Step 1 --- Start from the right-hand side.} Replace $F(\sigma)$ by its defining integral, $F(\sigma) = \int_0^\infty e^{-\sigma t} f(t)\,dt$, and write the result as a double integral:
\[ \color{workpurple} \int_s^\infty F(\sigma)\,d\sigma = \int_s^\infty\!\left(\int_0^\infty e^{-\sigma t} f(t)\,dt\right)d\sigma. \]
\textbf{Step 2 --- Exchange the order of integration.} Both variables are independent over the region $\sigma\in[s,\infty)$, $t\in[0,\infty)$, so (by Fubini, valid where the transform converges absolutely) we integrate in $\sigma$ first:
\[ \color{workpurple} \int_s^\infty F(\sigma)\,d\sigma = \int_0^\infty f(t)\left(\int_s^\infty e^{-\sigma t}\,d\sigma\right)dt. \]
\textbf{Step 3 --- Evaluate the inner $\sigma$-integral.} Treating $t>0$ as constant,
\[ \color{workpurple} \int_s^\infty e^{-\sigma t}\,d\sigma = \left[-\frac{1}{t}e^{-\sigma t}\right]_{\sigma=s}^{\sigma\to\infty} = 0 - \left(-\frac{1}{t}e^{-st}\right) = \frac{e^{-st}}{t}. \]
Substituting this back into the $t$-integral,
\[ \color{workpurple} \int_s^\infty F(\sigma)\,d\sigma = \int_0^\infty f(t)\,\frac{e^{-st}}{t}\,dt = \int_0^\infty e^{-st}\,\frac{f(t)}{t}\,dt. \]
\textbf{Step 4 --- Recognize the transform.} The last integral is, by definition, $\mathcal{L}\!\left\{\dfrac{f(t)}{t}\right\}$:
\[ \color{workpurple} \mathcal{L}\!\left\{\frac{f(t)}{t}\right\} = \int_s^\infty F(\sigma)\,d\sigma. \qquad\blacksquare \]
\textbf{Corollary (evaluation at $s=0$).} Setting $s=0$ makes $e^{-st}=1$ on the left, so $\mathcal{L}\{f/t\}\big|_{s=0} = \int_0^\infty \dfrac{f(t)}{t}\,dt$, giving
\[ \color{workpurple} \int_0^\infty \frac{f(t)}{t}\,dt = \int_0^\infty F(\sigma)\,d\sigma. \]
This is exactly the engine behind Solutions 1, 2, and 4.

\bigskip
\noindent\textbf{Solution 10 (Proof): Inversion of the Convolution Product.}\\[2pt]
\textbf{Goal.} Show $\mathcal{L}\{(f*g)(t)\} = F(s)\,G(s)$.

\smallskip
\textbf{Step 1 --- Write the product as a double integral.} Use a distinct dummy variable in each transform so they do not collide:
\[ \color{workpurple} F(s)\,G(s) = \left(\int_0^\infty e^{-s\tau} f(\tau)\,d\tau\right)\!\left(\int_0^\infty e^{-s\rho} g(\rho)\,d\rho\right) = \int_0^\infty\!\!\int_0^\infty e^{-s(\tau+\rho)} f(\tau)\,g(\rho)\,d\rho\,d\tau. \]
\textbf{Step 2 --- Substitute $t = \tau + \rho$ in the inner integral.} Hold $\tau$ fixed and let $t = \tau + \rho$, so $\rho = t - \tau$ and $d\rho = dt$. As $\rho$ runs from $0$ to $\infty$, the new variable $t$ runs from $\tau$ to $\infty$:
\[ \color{workpurple} F(s)\,G(s) = \int_0^\infty\!\!\int_{\tau}^\infty e^{-st} f(\tau)\,g(t-\tau)\,dt\,d\tau. \]
\textbf{Step 3 --- Swap the order of integration.} The region is $0 \le \tau \le t < \infty$. Reading it with $t$ outermost, $\tau$ ranges from $0$ to $t$:
\[ \color{workpurple} F(s)\,G(s) = \int_0^\infty e^{-st}\!\left(\int_0^t f(\tau)\,g(t-\tau)\,d\tau\right)dt. \]
\textbf{Step 4 --- Recognize the convolution.} The inner integral is precisely $(f*g)(t)$, so the outer integral is its Laplace transform:
\[ \color{workpurple} F(s)\,G(s) = \int_0^\infty e^{-st}\,(f*g)(t)\,dt = \mathcal{L}\{(f*g)(t)\}. \qquad\blacksquare \]
Read backwards, this says the \emph{product} $F(s)G(s)$ inverts to the \emph{convolution} $f*g$ --- the identity that turned every Volterra equation in this section into one line of algebra.

\begin{conceptbox}[Closing Takeaway]
The two properties proved above are the spine of the entire section. The \textcolor{accentorange}{integral property} ($\mathcal{L}\{f/t\} = \int_s^\infty F\,d\sigma$, evaluated at a point) lets a transform table \emph{evaluate} improper integrals that have no elementary antiderivative. The \textcolor{accentorange}{convolution theorem} ($\mathcal{L}\{f*g\} = F\,G$) lets the same table \emph{solve} for an unknown buried inside an integral. Whether the goal is a number or a function, the move is identical: cross the bridge into the $s$-domain, do easy algebra, and carry the answer back.
\end{conceptbox}

\end{document}
